% !TEX program = xelatex
\documentclass[11pt,a4paper]{article}
\newcommand{\DOCFOOT}{AMPERE POINT · Opis przypadku: wallbox + DLB-A1 · 26 września 2026}
% Wspólna preambuła raportu do producenta (xelatex)
\usepackage{fontspec}
% \ZHDOC zdefiniowane przed \input => cały dokument po chińsku (Noto Sans CJK SC + łamanie wierszy CJK)
\ifdefined\ZHDOC
  \setmainfont{Noto Sans CJK SC}[Scale=0.92]
  \setsansfont{Noto Sans CJK SC}[Scale=0.92]
  \XeTeXlinebreaklocale "zh"
  \XeTeXlinebreakskip = 0pt plus 1pt minus 0.1pt
\else
  \setmainfont{Lato}[Scale=0.95]
  \setsansfont{Lato}[Scale=0.95]
\fi
\setmonofont{DejaVu Sans Mono}[Scale=0.80]
\newfontfamily\cjkfont{Noto Sans CJK SC}[Scale=0.92]
\newfontfamily\latinfont{Lato}[Scale=0.95]
\newcommand{\pl}[1]{{\latinfont #1}}
\usepackage[a4paper,top=2.0cm,bottom=2.0cm,left=1.9cm,right=1.9cm]{geometry}
\usepackage{graphicx}
\usepackage[table]{xcolor}
\usepackage{booktabs,longtable,array,tabularx,enumitem,fancyhdr,caption}
\usepackage[most]{tcolorbox}
\usepackage[hidelinks,unicode]{hyperref}
\usepackage{needspace}
\usepackage{float}
\emergencystretch=2.5em
\setlength{\parindent}{0pt}\setlength{\parskip}{5pt}
\renewcommand{\arraystretch}{1.22}
\captionsetup{font=small,labelfont=bf,skip=3pt}
\setlist{nosep,leftmargin=*}
\setcounter{secnumdepth}{2}

\definecolor{apnavy}{HTML}{1A1A1A}
\definecolor{apink}{HTML}{1C2430}
\definecolor{apmuted}{HTML}{5B6675}
\definecolor{apline}{HTML}{E3E3E3}
\definecolor{apbg}{HTML}{F5F5F5}
\definecolor{apaccent}{HTML}{EC6434}
\definecolor{apaccentbg}{HTML}{FDEDE6}
\definecolor{apwarn}{HTML}{B9791B}
\definecolor{apwarnbg}{HTML}{FBF1DE}
\definecolor{apred}{HTML}{B91C1C}
\definecolor{apredbg}{HTML}{FEE2E2}
\definecolor{apgreen}{HTML}{1F8A4C}
\definecolor{apgreenbg}{HTML}{E7F5EC}
\definecolor{aphead}{HTML}{EDEDED}

\pagestyle{fancy}\fancyhf{}
\renewcommand{\headrulewidth}{0pt}
\fancyfoot[L]{\footnotesize\color{apmuted}\DOCFOOT}
\fancyfoot[R]{\footnotesize\color{apmuted}\thepage\,/\,\pageref{LastPage}}
\usepackage{lastpage}

\usepackage{titlesec}
\titleformat{\section}{\Large\bfseries\color{apnavy}}{\thesection}{0.6em}{}
\titlespacing*{\section}{0pt}{14pt}{6pt}
\titleformat{\subsection}{\large\bfseries\color{apnavy}}{\thesubsection}{0.6em}{}
\titlespacing*{\subsection}{0pt}{10pt}{4pt}

% pasek tytułowy
\newcommand{\apheader}[3]{%
\noindent\includegraphics[height=1.0cm]{logo_ap.png}\hfill
\raisebox{0.25cm}{\parbox[t]{10.5cm}{\raggedleft\footnotesize\color{apmuted}#3}}\\[0.25cm]
{\color{apaccent}\rule{\textwidth}{1.8pt}}\\[0.30cm]
{\bfseries\color{apnavy}\LARGE #1}\\[0.18cm]
{\large\color{apmuted}#2}\\[0.2cm]}

% ramki
\newtcolorbox{apbox}[1][]{enhanced,breakable,colback=apbg,colframe=apline,boxrule=0.6pt,arc=3pt,left=7pt,right=7pt,top=5pt,bottom=5pt,fonttitle=\bfseries,coltitle=apnavy,colbacktitle=aphead,title={#1}}
\newtcolorbox{apwarnb}[1][]{enhanced,breakable,colback=apwarnbg,colframe=apwarn,boxrule=0.6pt,arc=3pt,left=7pt,right=7pt,top=5pt,bottom=5pt,fonttitle=\bfseries,title={#1}}
\newtcolorbox{apredb}[1][]{enhanced,breakable,colback=apredbg,colframe=apred,boxrule=0.6pt,arc=3pt,left=7pt,right=7pt,top=5pt,bottom=5pt,fonttitle=\bfseries,title={#1}}
\newtcolorbox{apgreenb}[1][]{enhanced,breakable,colback=apgreenbg,colframe=apgreen,boxrule=0.6pt,arc=3pt,left=7pt,right=7pt,top=5pt,bottom=5pt,fonttitle=\bfseries,title={#1}}

% komórka ze zrzutem ekranu: #1 szerokość, #2 plik, #3 podpis
\newcommand{\shot}[3]{%
\begin{minipage}[t]{#1}
  \vspace{0pt}\centering
  \includegraphics[width=\linewidth,height=6.2cm,keepaspectratio]{#2}\par\vspace{3pt}
  {\footnotesize\color{apink}#3\par}
\end{minipage}}


% mniejsza komórka ze zdjęciem (np. gniazda aut)
\newcommand{\shotsm}[3]{%
\begin{minipage}[t]{#1}
  \vspace{0pt}\centering
  \includegraphics[width=\linewidth,height=4.6cm,keepaspectratio]{#2}\par\vspace{3pt}
  {\footnotesize\color{apink}#3\par}
\end{minipage}}

\newcommand{\kroknum}[1]{{\colorbox{apaccent}{\textcolor{white}{\bfseries\small\,#1\,}}}}
\newcommand{\dpn}[1]{{\cjkfont #1}}
\newcommand{\Fact}{\textcolor{apgreen}{\textbf{[fact]}}}
\newcommand{\Hyp}{\textcolor{apwarn}{\textbf{[hypothesis]}}}

\usepackage{polyglossia}\setdefaultlanguage{polish}
\renewcommand{\Hyp}{\textcolor{apwarn}{\textbf{[hipoteza]}}}
\renewcommand{\figurename}{Rys.}
\hypersetup{pdftitle={Opis przypadku: wallbox z kontrolerem DLB-A1},pdfauthor={AMPERE POINT Serwis}}
\begin{document}

\apheader{Opis przypadku: wallbox z kontrolerem DLB-A1}%
{Prąd ładowania przekracza nastawę wallboxa, zadziałanie ochrony przed przeciążeniem, brak wznowienia ładowania po pauzie DLB}%
{AMPERE POINT — Serwis · zgłoszenie do producenta\\26 września 2026\\Sprawa: wallbox \texttt{bfe26d78713489f8a6ulsj} + DLB-A1}

\section{Streszczenie}

{\sloppy Wallbox AmperePoint 11~kW, trójfazowy (ID urządzenia Tuya \texttt{bfe26d78713489f8a6ulsj}, firmware \texttt{(V8.0.7)F1.3.6}) pracował u klienta poprawnie. Po dołożeniu do instalacji kontrolera dynamicznego zarządzania mocą DLB-A1 pojawiły się dwa problemy. 26 września 2026 instalatorzy stwierdzili oba problemy na miejscu i udokumentowali je nagraniem wideo. Te same zdarzenia widać w logu urządzenia na platformie Tuya. Wszystkie fakty w tym raporcie pochodzą z tej dokumentacji wideo i z odpowiadających jej wpisów w logu urządzenia.\par}

\begin{apredb}[Problem 1 — absolutny priorytet: w trybie DLB wallbox nie respektuje własnej nastawy prądu]
Przy nastawie wallboxa \textbf{16~A} i kontrolerze DLB ustawionym na \textbf{18~A} wallbox ładował prądem \textbf{16,9~A} (7,3~kW). Po obniżeniu nastawy wallboxa do \textbf{9~A} w trakcie ładowania wallbox przez kilka sekund ładował dalej prądem \textbf{16,6~A / 7,1~kW}, po czym wyświetlił \emph{„Ochrona przed przeciążeniem prądu — Automatyczne przełączenie na ładowanie niskim prądem”}; zamiast przełączyć się na niski prąd, przestał ładować w ogóle. \textbf{Problem 1 wywołał awarię ładowarki pokładowej w samochodzie trójfazowym: samochód przestał reagować i do przywrócenia go do pracy potrzebny był mechanik.}
\end{apredb}

\begin{apwarnb}[Problem 2: po pauzie DLB wallbox nie wznawia ładowania]
Gdy dodatkowe obciążenie domu sprawia, że DLB obniża prąd wallboxa poniżej 6~A, wallbox przestaje ładować — poprawnie. Po odłączeniu obciążenia ładowanie \textbf{nie} wznawia się, nawet po 5 minutach. Wznawia się dopiero po odłączeniu i ponownym podłączeniu samochodu.
\end{apwarnb}

Przed dołożeniem DLB wallbox pracował poprawnie na tej samej instalacji i tym samym gnieździe, więc instalacja jest wykluczona. Testy wykonano na dwóch samochodach — jednym z dwufazową i jednym z trójfazową ładowarką pokładową — więc samochód jest wykluczony. Użyto dwóch niezależnych zestawów kontrolerów DLB z akcesoriami i problemy wystąpiły na obu. To jest zatem problem zestawu wallbox / DLB. Prosimy o analizę i poprawione oprogramowanie. Poniżej dowody, fragmenty logu i konkretne pytania.

\section{Sprzęt}

\begin{longtable}{>{\raggedright\arraybackslash}p{3.2cm}>{\raggedright\arraybackslash}p{12.9cm}}
\toprule
\rowcolor{aphead}\textbf{Element} & \textbf{Szczegóły} \\
\midrule\endhead
\bottomrule\endfoot
Wallbox & Wallbox AmperePoint, 11~kW, trójfazowy, maksymalnie 16~A. ID urządzenia Tuya \texttt{bfe26d78713489f8a6ulsj}, ID produktu \texttt{gbmxngploofmhbjc}, firmware \texttt{(V8.0.7)F1.3.6}, informacja o urządzeniu \emph{„Type~B, AC~30mA + DC~6mA”}, wariant produktu 3 (dane raportowane przez urządzenie w logu Tuya). \\
Kontroler DLB & DLB-A1. Użyto \textbf{dwóch niezależnych zestawów} (kontroler z akcesoriami) — problemy wystąpiły na obu. \\
Samochody & \textbf{Samochód A}: dwufazowa ładowarka pokładowa (ładuje tylko na L1 i L2) — Problem 1. \textbf{Samochód B}: trójfazowa ładowarka pokładowa — Problem 2; jego ładowarka pokładowa uległa awarii wskutek Problemu 1. Gniazda ładowania obu samochodów: Rys.~\ref{fig:cars}. \\
\end{longtable}

\begin{apbox}[Co wykluczono]
\begin{itemize}
\item \textbf{Instalacja:} przed dołożeniem DLB wallbox ładował poprawnie na tej samej instalacji i tym samym gnieździe.
\item \textbf{Wadliwy wallbox:} ten sam wallbox bez DLB działa poprawnie.
\item \textbf{Samochód:} przetestowano na dwóch różnych samochodach (dwufazowa i trójfazowa ładowarka pokładowa).
\item \textbf{Wadliwy egzemplarz DLB albo przekładnik prądowy:} użyto dwóch pełnych, niezależnych zestawów DLB-A1 (kontroler ze wszystkimi akcesoriami, w tym przekładnikami) — problemy wystąpiły na obu.
\end{itemize}
\end{apbox}

\section{Problem 1 — prąd powyżej nastawy wallboxa, zadziałanie ochrony przed przeciążeniem, zatrzymanie ładowania}

\subsection{Udokumentowane na nagraniu (26 września 2026, 12:27–12:32 czasu lokalnego)}

Samochód A (dwufazowy). Ustawienia: kontroler DLB \textbf{18~A} (wyświetlacz DLB \emph{018}, Rys.~\ref{fig:p1}), nastawa wallboxa \textbf{16~A}.

\begin{enumerate}[leftmargin=2.2em]
\item \kroknum{1} Trwa ładowanie. Wyświetlacz wallboxa: \textbf{DLB 16,9~A}, 7,3~kW, 225~V, nastawa 16~A. Pomimo nastawy 16~A na wallboxie DLB ogranicza prąd do 16,9~A — nawet z uwzględnieniem pewnej tolerancji jest to powyżej nastawy.
\item \kroknum{2} Nastawa prądu ładowania zostaje obniżona na wallboxie z 16~A do \textbf{9~A} w trakcie ładowania. Wyświetlacz: DLB 16,9~A, 7,3~kW, nastawa 9~A.
\item \kroknum{3} Wallbox nie zmniejsza mocy ładowania: przez kilka sekund ładuje dalej prądem \textbf{16,6~A / 7,1~kW} przy nastawie 9~A.
\item \kroknum{4} Po kilku sekundach pojawia się komunikat ochrony przed przeciążeniem: \emph{„Ochrona przed przeciążeniem prądu — Automatyczne przełączenie na ładowanie niskim prądem”}. Wyświetlacz: DLB 0,0~A, 226~V, nastawa już \textbf{8~A}. Wallbox nie przełączył się na ładowanie niskim prądem — przestał ładować w ogóle.
\end{enumerate}

\begin{figure}[H]
\centering
\shot{0.19\textwidth}{p1_step1.png}{\kroknum{1} 16,9~A, 7,3~kW, nastawa 16~A}\hfill
\shot{0.19\textwidth}{p1_step2.png}{\kroknum{2} nastawa obniżona do 9~A — nadal 16,9~A}\hfill
\shot{0.19\textwidth}{p1_step3.png}{\kroknum{3} kilka sekund później: 16,6~A, 7,1~kW przy 9~A}\hfill
\shot{0.19\textwidth}{p1_step4.png}{\kroknum{4} ochrona przed przeciążeniem, nastawa 8~A, ładowanie zatrzymane}\hfill
\shot{0.19\textwidth}{p1_dlb018_rot.png}{Wyświetlacz DLB: SEt = 018 (18~A); klatka obrócona o 180° dla czytelności}
\caption{Problem 1 — klatki z nagrania instalatorów.}
\label{fig:p1}
\end{figure}

\subsection{Te same minuty w logu urządzenia (platforma Tuya)}

Czasy lokalne (UTC+2). Nazwy punktów danych cytowane tak, jak występują w logu Tuya.

\begin{longtable}{>{\raggedright\arraybackslash}p{2.0cm}>{\raggedright\arraybackslash}p{14.1cm}}
\toprule
\rowcolor{aphead}\textbf{Czas} & \textbf{Zdarzenie} \\
\midrule\endhead
\bottomrule\endfoot
12:27:14 & Wallbox online po włączeniu zasilania; \dpn{设置充电电流} (nastawa prądu) = 16; 12:27:20 \dpn{工作状态} (stan pracy) = 300 (ładowanie) \\
12:29:11 – 12:31:10 & \dpn{指标信息} (pomiary): \textbf{L1 17,0~A, L2 15,3~A, L3 0~A, 7,3~kW} — stabilnie \\
12:31:57 – 12:31:59 & \dpn{设置充电电流} zgłaszane przez urządzenie (nie przez aplikację): 6 → 7 → 8 → \textbf{9}; stan pracy nadal 300 (ładowanie) \\
12:32:09 & \dpn{报警信息} (alarm) \texttt{\{"t":"2026-09-26 12:32:05","v":400\}}; \dpn{工作状态} = \textbf{400}; \dpn{调试用工作状态} = \dpn{发生故障} (błąd); \dpn{设置充电电流} = \textbf{8}; pomiary 0~A \\
12:32:30 & \dpn{充电记录} (rekord sesji): 12:27–12:32, 287~s \\
\end{longtable}

Log zgadza się z nagraniem: 7,3~kW przy nastawie 16~A, nastawa obniżona do 9~A na urządzeniu w trakcie ładowania, brak obniżenia prądu, potem alarm 400 ok. 10~s po ostatniej zmianie nastawy, nastawa obniżona do 8~A i ładowanie zatrzymane. Pomiar prądu przez wallbox w logu (17,0~A na L1) zgadza się z wartością w polu DLB na wyświetlaczu (16,9~A).

\begin{apredb}[Skutek]
Problem 1 wywołał awarię ładowarki pokładowej w samochodzie trójfazowym. Samochód przestał reagować i do przywrócenia go do życia potrzebny był mechanik. Naprawa Problemu 1 jest zatem absolutnym priorytetem.
\end{apredb}

\section{Problem 2 — brak automatycznego wznowienia po pauzie DLB}

\subsection{Udokumentowane na nagraniu (26 września 2026, 12:36–12:44 czasu lokalnego)}

Samochód B (trójfazowy). Ustawienia: kontroler DLB \textbf{celowo ustawiony na mniejszy prąd}, aby wywołać przypadek wystąpienia błędu, co daje mniejszy maksymalny prąd wallboxa; nastawa wallboxa 16~A.

\begin{enumerate}[leftmargin=2.2em]
\item \kroknum{1} Ładowanie przy \textbf{DLB 10,8~A}, 7,4~kW, 226~V, nastawa 16~A.
\item \kroknum{2} Podłączony zostaje czajnik, aby zasymulować dodatkowe obciążenie instalacji. DLB działa tak, jak powinien, i obniża prąd wallboxa; prąd spada poniżej 6~A i wallbox \textbf{poprawnie przestaje ładować}. Wyświetlacz: DLB 0,0~A, 0,0~kW, 228~V, nastawa 16~A, samochód nadal podłączony.
\item \kroknum{3} Czajnik zostaje odłączony. \textbf{Wallbox nie zaczyna ładować ponownie.} Instalatorzy odczekali 5 minut — nadal brak ładowania.
\item \kroknum{4} Ładowanie wznawia się dopiero po \textbf{odłączeniu i ponownym podłączeniu samochodu}.
\end{enumerate}

\begin{figure}[H]
\centering
\shot{0.26\textwidth}{p2_step1.png}{\kroknum{1} ładowanie przy DLB 10,8~A, 7,4~kW, nastawa 16~A}\hspace{1.5cm}
\shot{0.26\textwidth}{p2_step2.png}{\kroknum{2} zatrzymanie po podłączeniu czajnika: DLB 0,0~A, 0,0~kW, samochód podłączony — i tak zostaje}
\caption{Problem 2 — klatki z nagrania instalatorów.}
\label{fig:p2}
\end{figure}

\subsection{Te same minuty w logu urządzenia (platforma Tuya)}

\begin{longtable}{>{\raggedright\arraybackslash}p{2.0cm}>{\raggedright\arraybackslash}p{14.1cm}}
\toprule
\rowcolor{aphead}\textbf{Czas} & \textbf{Zdarzenie} \\
\midrule\endhead
\bottomrule\endfoot
12:35:51 & Wallbox online po restarcie zasilania; trwa ładowanie (pomiary prądu od 12:35:54) \\
12:36:13 – 12:37:34 & \dpn{指标信息}: \textbf{10,5–10,8 / 10,7–11,0 / 11,1–11,6~A, 7,4~kW}, 225–226~V — stabilnie \\
12:37:53 & \dpn{工作状态} = \textbf{200} (samochód podłączony, nie ładuje) — pauza \\
12:37:53 – 12:44:21 & Brak zgłoszenia stanu ładowania i prądu; wallbox pozostaje w stanie 200 \\
12:44:21 & Wallbox przechodzi w tryb offline (\emph{keepalive\_timeout}) \\
\end{longtable}

Instrukcja DLB-A1 Quick User Manual (sekcja 6.2, \emph{Paused}) mówi: \emph{„Recovery is automatic when headroom returns.”} (wznowienie następuje automatycznie po powrocie zapasu mocy). To nie następuje.

\section{Zachowanie zaobserwowane a udokumentowane}

\begin{longtable}{>{\raggedright\arraybackslash}p{5.0cm}>{\raggedright\arraybackslash}p{5.1cm}>{\raggedright\arraybackslash}p{5.8cm}}
\toprule
\rowcolor{aphead}\textbf{Funkcja} & \textbf{Udokumentowane (DLB-A1 Quick User Manual, Rev.~20251109) albo wyświetlane} & \textbf{Zaobserwowane} \\
\midrule\endhead
\bottomrule\endfoot
Górna granica prądu ładowania & „each wall box is capped by its own requested current” (każdy wallbox jest ograniczony do własnego żądanego prądu) & Prąd przekracza nastawę wallboxa: 16,9~A przy nastawie 16~A; 16,6~A przy nastawie 9~A \\
Limit użyteczny & „90\,\% of the Configured Supply Limit (10\,\% buffer)” (90\,\% skonfigurowanego limitu, 10\,\% bufora) & Przy DLB ustawionym na 18~A wallbox ładował prądem 16,9~A, powyżej 90\,\% $\times$ 18~A = 16,2~A \\
Redukcja przy spadku zapasu, pauza poniżej 6~A & redukuje, utrzymuje $\geq$6~A gdzie to możliwe, inaczej pauza & Działa (Problem 2, krok 2) \\
Wznowienie po pauzie & „Recovery is automatic when headroom returns” (wznowienie automatyczne po powrocie zapasu) & Brak wznowienia, nawet po 5~min; konieczne ponowne podłączenie samochodu \\
Komunikat ochrony przed przeciążeniem & Komunikat wallboxa: „Automatyczne przełączenie na ładowanie niskim prądem” & Ładowanie staje całkowicie; nastawa obniżona z 9~A na 8~A \\
\end{longtable}

\section{Nasze odczytanie dowodów}

Poniżej nasze odczytanie dowodów; prosimy o jego potwierdzenie albo skorygowanie.

\begin{enumerate}
\item W trybie DLB pozwolenie prądowe sygnalizowane samochodowi na Control Pilot wynika, jak się wydaje, wyłącznie z przydziału DLB i \textbf{nie jest ograniczane własną nastawą wallboxa}: 16,9~A przy nastawie 16~A, 16,6~A przy nastawie 9~A. Przy DLB ustawionym na 18~A pozwolenie było też powyżej 90\,\% tej nastawy, co wskazuje, że 10-procentowy bufor bezpieczeństwa nie jest stosowany.
\item Gdy nastawa jest obniżana w trakcie ładowania, nowa wartość jest, jak się wydaje, od razu stosowana w ochronie przed przeciążeniem, a pozwolenie na CP zostaje na wartości z DLB. Samochód nigdy nie dostaje polecenia obniżenia prądu, a ok. 10~s później ochrona zadziała, obniża nastawę o 1~A i zatrzymuje ładowanie.
\item Problem 2 jest osobny: po pauzie spowodowanej przydziałem poniżej 6~A wallbox (albo DLB) nie wznawia ładowania po odłączeniu dodatkowego obciążenia.
\end{enumerate}

Oczekiwane zachowanie: pozwolenie na CP = min(nastawa wallboxa, przydział DLB); zmiana nastawy w trakcie ładowania najpierw obniża wypełnienie PWM i daje samochodowi standardowy czas na dostosowanie się, zanim jakakolwiek ochrona zostanie oceniona; automatyczne wznowienie po pauzie.

\section{Prośby i pytania}

\begin{enumerate}
\item \textbf{Przede wszystkim: prosimy o naprawienie obu problemów, tak aby nie występowały w kolejnych egzemplarzach.} Prosimy o poprawione oprogramowanie wallboxa (obecnie \texttt{(V8.0.7)F1.3.6}) i/lub DLB-A1 oraz o informację, których wersji oprogramowania i partii produkcyjnych to dotyczy, abyśmy mogli poinformować pozostałych klientów.
\item Czy w trybie DLB pozwolenie na CP jest liczone jako \textbf{min(nastawa wallboxa, przydział DLB)}? Nasze obserwacje wskazują, że nie.
\item Jaki wzór przydziału faktycznie realizuje DLB-A1? Instrukcja podaje 90\,\% $\times$ skonfigurowany limit zasilania minus najwyższy prąd fazowy zmierzony przez CT; przy DLB ustawionym na 18~A zaobserwowaliśmy 16,9~A.
\item Jaki jest warunek zadziałania ochrony przed przeciążeniem (alarm 400) i co dokładnie robi „automatyczne przełączenie na ładowanie niskim prądem”? Zaobserwowaliśmy: ładowanie staje całkowicie, a nastawa jest obniżana o 1~A.
\item Czy przy zmianie nastawy w trakcie ładowania wallbox najpierw obniża pozwolenie na CP i czeka, aż samochód się dostosuje, zanim oceni ochronę?
\item Dlaczego wallbox nie wznawia ładowania po pauzie DLB, gdy dodatkowe obciążenie zostaje odłączone? Czy ponowne podłączenie samochodu jest wymagane z założenia?
\item Na potrzeby diagnostyki: czy przydział DLB (prąd docelowy wysyłany do wallboxa) mógłby być raportowany jako punkt danych w logu Tuya? Obecnie log nie zawiera żadnego punktu danych związanego z DLB.
\end{enumerate}

\needspace{6\baselineskip}
\begin{apgreenb}[Rozwiązanie tymczasowe, które przekazujemy klientom]
Do czasu poprawionego oprogramowania: praca wallboxa bez DLB — bez niego wallbox działa poprawnie.
\end{apgreenb}

\section{Załączniki i dostępne materiały}

\begin{itemize}
\item Klatki z nagrania instalatorów z 26 września 2026 (Problem 1 i Problem 2): Rys.~\ref{fig:p1} i Rys.~\ref{fig:p2}.
\item Zdjęcia gniazd ładowania obu samochodów (Rys.~\ref{fig:cars}).
\item Log urządzenia \texttt{bfe26d78713489f8a6ulsj} z 26 września 2026 na platformie Tuya — cytowane wyżej zdarzenia można tam odnaleźć po tym ID urządzenia. Wszystkie czasy w tym raporcie są lokalne (UTC+2); sam log jest zapisany w UTC.
\end{itemize}

\begin{figure}[H]
\centering
\shotsm{0.24\textwidth}{car2ph.png}{Samochód A — dwufazowa ładowarka pokładowa}\hspace{1.2cm}
\shotsm{0.24\textwidth}{car3ph.png}{Samochód B — trójfazowa ładowarka pokładowa}
\caption{Samochody użyte w testach.}
\label{fig:cars}
\end{figure}

{\color{apmuted}\small AMPERE POINT — Serwis · www.amperepoint.pl · +48 666 612 044 · raport sporządzony 26 września 2026 na podstawie nagrania instalatorów i logu urządzenia Tuya.}

\end{document}
